\subsection{Методология сравнения групп признаков}

Экспериментальная методология строится как сравнение групп признаков при едином протоколе обучения и оценки. Это означает, что различия в результатах должны интерпретироваться прежде всего как различия в информативности входных признаков, а не как следствие разных условий эксперимента.

Основные группы признаков сравниваются в одинаковой задаче бинарной классификации. Для всех вариантов используется один и тот же target:

\[
y_t = \mathbf{1}\{Close_{t+1} > Close_t\}.
\]

Разбиение данных выполняется хронологически. Сначала модель обучается на train subset, затем используется validation subset, а итоговая оценка проводится на test subset. Одна и та же дата не может попасть одновременно в разные подвыборки. Это важно для финансовых временных рядов, поскольку случайное перемешивание наблюдений может привести к утечке информации между близкими датами.

Нормализация рыночных признаков выполняется с помощью \texttt{StandardScaler}. Scaler обучается только на train subset, после чего то же преобразование применяется к validation и test. Такой порядок нужен для предотвращения утечки данных из будущих периодов в preprocessing.

Метрики качества рассчитываются на test subset. Для оценки используются accuracy, balanced accuracy, precision, recall, F1-score, ROC-AUC и confusion matrix. Accuracy показывает общую долю правильных предсказаний, однако при несбалансированном распределении классов она может быть недостаточной. Поэтому дополнительно анализируются balanced accuracy и ROC-AUC, позволяющие оценить способность модели различать оба класса.

Методология сравнения строится вокруг следующих вопросов:

\begin{enumerate}
    \item содержат ли FinBERT-признаки самостоятельный сигнал относительно majority baseline;
    \item дают ли рыночные признаки прогнозный сигнал относительно majority baseline;
    \item улучшают ли FinBERT-признаки качество прогноза при добавлении к рыночным признакам;
    \item усиливает ли объединение текстовой и рыночной информации результат относительно отдельных групп признаков;
    \item помогает ли более сложная последовательная модель извлекать дополнительный сигнал из объединённых признаков.
\end{enumerate}

Такой дизайн соответствует feature-centric логике исследования: модели используются как инструмент проверки признаков, а не как самостоятельный объект оптимизации.
